\iffalse
(Примечание автора: А чего выводы находятся в введении? Пусть в конце будет и в отдельному документе, что-то не так как будто.)
\fi

\clearpage
\section{Выводы}
\label{sec:ch04_conclusions}
\conceptsource{Синтез статей 1--4 по согласованному автором направлению
R4-17. Архитектурные возможности, проверенные эффекты и развитие
исследовательских задач описаны совместно.}

В главе рассмотрено развитие нейросетевой обработки
электрических сигналов от оценки события в окне
к разметке отдельных участков переходного процесса.
Для включения электротехнических зависимостей разработаны
физические ветви с произведениями и отношениями величин,
комплексными представлениями и фазовым управлением
амплитудными признаками. Эти идеи объединены с KAN
и вниманием в архитектуре Physical KAN-Transformer.
Внимание связывает временные участки, а физические
блоки задают обработку отношений между электрическими величинами.
\iffalse
(Примечание автора: «Показал использование»? Странно очень звучит, перепиши. Да и производные признаки здесь и в других местах не совсем понятно, путается с понятием "производной сигнала - dx/dt", может чуть другой терминн придумать получится?)
(Примечание автора: Ок, закрываем, прост для информации - добавил к предыдущему абзацу.)
\fi
\authorreviewnote{AA131}
Вклад физических преобразований проверен отключением отдельных
ветвей обученной PhysicsKAN. Отключение произведений
и исходных входов вызывает наибольшее снижение качества
в рассмотренном опыте. Небольшое обнуление слабых связей
сохраняет общий результат, тогда как более сильное сокращение
увеличивает ошибки. Тем самым проверена полезность
сочетания исходных данных с физически рассчитанными признаками.

\iffalse
(Примечание автора: Как показало? Какое сопоставление? Перепиши первую фразу. Да и не только...
Да и весь абзац плоховат. Ты здесь всего лишь хотела сказать о том, что хороши представления через модули и фазы исходных сигналов. А так же неплохо всё в последовательностях представлять. А вторая идея о том, что модели для итоговой задачи можно бы подбирать как комбинацию, ибо меняется сложность – вычислительная скорость и прочие, да и этот момент можно бы отдельно. Мол мы не утверждаем, что пихать надо всегда сверх тяжёлую модель, а надо бы рационально выбирать, и тот же РНМ мы будет разреженным брать из трансформера физичного, прост будет точку предсказания ставить в конце
Это про идущее ниже.)
(Примечание автора: Ок, закрываем вопрос, чуть переписал сам и сойдёт.)
\fi
\authorreviewnote{AA132}
В исследовании входных представлений высокие оценки получены
при ипользовании спектральных данных (модуль + фаза). Такой вход сохраняет
амплитуды и угловые соотношения электрических величин.
Последовательность спектральных состояний позволяет учитывать
изменение этих величин в переходном процессе.

Число входных состояний и сложность сети стоит далее выбирать опираясь
на желаемый результатом. Подробная разметка событий требует
несколько временных участков и стоит использовать "объёмные" по количеству входных и выходных данных модели.
При определении же состояния РНМ и большинства других органов РЗА и БАВР
считаем корректным применение разреженной спектральная последовательность,
а при необходимости и сжатой до двух значений, с одним ответом в конце доступного интервала. 

\iffalse
(Примечание автора: Тоже проверь и перепиши абзац ниже, ибо ты и закапываешься в слова… И лучшей вроде бы другая модель оказалась, как раз физичная КАН, и главная мысль тут в том, что обычный трансформер может и хорошо выучился на идеальных данных, но стал местами себя хуже вести на реальном. И не до конца понятна тут вот эта позонность, надо чуть проще про факт определения, про всякие зоны уже там локально пишется в случае чего.)
\fi
\authorreviewnote{AA133}
Проверка ОЗЗ выявила различие результатов на расчётных
и реальных процессах. На моделируемых сигналах наибольший
оконный F1 получен базовой спектральной моделью.
На реальных записях Physical KAN-Transformer обнаружила
94,81\% подтверждённых ОЗЗ, а базовая спектральная модель
92,64\%. При этом доля ложных обнаружений в отрицательной
группе составила соответственно 63,65 и 55,68\%.
Таким образом, обнаружение повреждения рассматривается
вместе с реакцией на другие переходные процессы.
Проверка на эксплуатационных осциллограммах дополняет
сравнение на идеализированных моделях повреждений.

Применение физической трансформерной обработки
к насыщению ТТ расширяет исследование на качество
измерительного тракта. Модель распознаёт участки искажения
вторичного тока по предоставленным расчётным меткам,
а реальные примеры позволяют рассмотреть изменение выхода
на зарегистрированных процессах.

Полученные результаты связывают развитие архитектур,
выбор данных и проверку на прикладных задачах.
Дальнейшие широкие вычислительные исследования могут уточнить
сочетания физических блоков, полезность каждого из них 
и актуальную сложности моделей для разных типов процессов.

В следующей главе рассматривается
общая методика проверки и модернизации органов РЗА,
включая адаптивные решения. А РНМ служит предметным
примером применения этой методики в контексте улучшения работы БАВР
с оценкой направления, применимости и поведения решения во времени.
\iffalse
(Примечание автора: Ну… В следующей главе не только модернизация БАВР, можно вопрос чуть более глобально ставить, о методике проверки разработанных органов РЗА, которые заявляются как прям адаптивные под любой режим, ну и РНМ как раз является им.)
(Примечание автора: Ок, закрывай вопрос)
\fi
\authorreviewnote{AA134}

